\subsection{Condiciones y supuestos de diseño}
\label{sec:contradicciones}

La lógica adopta seis instalaciones en total, de las cuales cinco son sitios logísticos con servicio local y la sexta es la casa matriz de Talca, contigua al CD principal y sin cómputo propio (Caso 02, cap.~2.3). La parametrización de sitios evita codificar el número en cada módulo.

Los cortes frecuentes de dos horas descritos en la operación son el escenario ordinario. Los límites de diseño son más exigentes: un turno de 14 horas en terreno y 24 horas en cada centro de distribución. No se reduce la autonomía porque el incidente más común sea más corto.

El relevo de turnos sin IdP, la autorización tributaria de salida y la pérdida de sitio se verifican con AL-OFF-01, AL-DTE-01 y AL-DR-01 del Anexo 4.1-M. Los objetivos exigidos siguen siendo RTO $\leq$ 4 horas y RPO $\leq$ 15 minutos para servicios críticos. La retención en el sitio no acredita recuperación ante su pérdida simultánea con la WAN. El registro de cierre distingue diseño, prueba ejecutada y aprobación del CLIENTE; ninguno sustituye a los otros.

\clearpage

El Anexo 4.1-V reúne las condiciones de aceptación y las decisiones del CLIENTE. La continuidad tributaria de las 96 salidas, la protección ante pérdida del sitio bajo aislamiento total y la revisión humana del consolidado requieren evidencia o aprobación identificable antes de declarar cierre; la propuesta no rebaja los requisitos por omisión.
